% =====================================================================
%  Chapter 18 — The ten walls
% =====================================================================
\chapter{The Ten Walls}
\label{ch:walls}

\begin{record}
\textbf{Source.} This chapter reports \texttt{sec:walls-s47}, the walls
document written in \Sn{47} and amended in \Sn{48}. It occupies lines
\(10493\)--\(11157\) of the proof document as it now stands.\footnote{The
commissioning brief for this chapter, and \texttt{PROGRESS.md}, both give the
range as \(9716\)--\(10348\). Those were the line numbers at \Sn{47}; \Sn{48}
added \(847\) lines to the document, of which \(777\) fall ahead of this
section --- and it also rewrote parts (c) and (d) of Wall~2 inside it. The
label \texttt{sec:walls-s47} is what should be cited, and the line numbers are
secondary for exactly this reason.}

\smallskip
\textbf{What a walls document is for.} It is a record, not an argument. Its
stated purpose is ``so that a later session does not repeat a step already
proved impossible, and so that the material exists in citable form outside the
session logs''. Its own scope remark is explicit that it contains no new
mathematics, that no statement in it is cited as a hypothesis anywhere else,
and that it may be deleted without altering any proof.

\smallskip
\textbf{The template.} Every wall carries four parts, and this chapter keeps
them: \emph{route} (what was attempted, and what it would have delivered),
\emph{obstruction} (the statement that stops it, by label),
\emph{classification} (exactly one label), and \emph{consequence} (what the
wall forbids and what it leaves permitted). The classifications are the
document's own and are reproduced without softening --- in particular Wall~1
is \OPEN{}, not \IMPOSSIBLE{}, and Wall~2 is \WITHDRAWN{}.
\end{record}

\section{Wall 0 --- the \texorpdfstring{\(H^1\)}{H1} Prize claim of \S20--\S22}
\label{sec:wall0}

\textbf{Route.} A vorticity scale-recursion iterated down scales to give the
uniform enstrophy cascade bound \(Z(r,z_0)\le D\,r^{2/3}\) for all
\(u_0\in H^1(\T^3)\), from which Claim~A and ``the Prize for \(u_0\in H^1\)''
were drawn.

\textbf{Obstruction.} \texttt{sec:audit-s31} (line~5218) records \emph{eight}
independent defects, not one. They are set out in full in
Chapter~\ref{ch:audit}; in one line each: (D1) the induction requires \(D\)
small and its base case requires \(D\) large; (D2) an invalid H\"older triple,
\(\tfrac13+\tfrac12+\tfrac12>1\); (D3) a reversed Jensen inequality; (D4)
regularity circularity plus global/local norm mixing; (D5) the Gr\"onwall
theorem assumes the Beale--Kato--Majda criterion, i.e.\ an equivalent of its
conclusion; (D6) the two-regime dichotomy needs a viscosity larger than
universal constants; (D7) the Regime~II ODE conclusion is false, since
\(\dot M=CM^{5/4}\) blows up in finite time; (D8) unproved coercivity,
\emph{together with} the false bound \(\norm{\nabla u}_{L^\infty}\le CM\).

\textbf{Classification.} \WITHDRAWN{}. Global regularity for general
\(u_0\in H^1(\T^3)\) is \textbf{open} in the document.

\textbf{Consequence.} \S20 remains usable only as a \emph{conditional
small-data} statement, and only after D2--D4 are repaired and the
superlinearity exponent re-derived. The audit's own survivors list is
reproduced in Chapter~\ref{ch:inheritance}. And, in the wall's closing clause:
``Defect~D8 is the first appearance of Wall~1 below; it was already flagged in
the audit and is the reason the logarithm is not a later regression.''

\section{Wall 1 --- the Calder\'on--Zygmund logarithm}
\label{sec:wall1}

\textbf{Route.} Every enstrophy-budget estimate from \Sn{33} to \Sn{44}
controls the vortex stretching by a sup-norm of \(\nabla u\), recovered from
\(\omega\) by Biot--Savart. Closing any of them requires that sup-norm to cost
\(O(M)\).

\textbf{Obstruction.} It costs \(M\Lfac\). The factor
\(\Lfac=\log(e+\norm{u}_{H^3}/M)\) enters at
\texttt{lem:local-enstrophy-typeI} (line~5701, definition on line~5712)
through \(\norm{\nabla u}_{L^\infty}\lesssim M\Lfac\), the \(L^\infty\)
endpoint failure of Biot--Savart. Every budget downstream inherits it, and it
appears at three distinct closure attempts:

\begin{enumerate}[label=\textup{(\roman*)},leftmargin=2.4em]
\item \emph{Crude \(K\)-absorption}, \texttt{prop:crude-K} (line~5903): the
  coupling term re-enters at top order multiplied by \(\Lfac\) --- ``the crude
  route misses closing the loop by exactly one logarithm''.
\item \emph{The annulus core}, \texttt{rem:s41-shortfall} (line~8428):
  absorption requires \(C\Lfac\le\delta\nu\), which fails for fixed \(\nu\) as
  \(\Lfac\to\infty\); the shortfall is ``exactly the same single logarithm''.
\item \emph{The corrected band route}, \texttt{prop:s44-linfty-route}
  (line~9242): both remaining conjectures close if and only if the correlation
  constant \emph{decays} like \(C_1/\Lfac\); boundedness is not enough, because
  the hypothesis is used precisely to cancel the \(\Lfac\) the budget carries.
\end{enumerate}

\textbf{Classification.} \OPEN{}, and structural. \emph{No proof in the
document says the logarithm cannot be removed.} What the three items establish
is that it is not a bookkeeping artefact of any one estimate: it is paid once,
at the point where a sup-norm of a Calder\'on--Zygmund image is taken, and then
inherited.

\textbf{Consequence.} \emph{Forbidden:} presenting (i)--(iii) as separate
obstacles to be attacked separately --- they are one obstacle seen three times.
\emph{Forbidden:} closing any of them with a \emph{boundedness} hypothesis on a
correlation constant; the required hypothesis is decay. \emph{Permitted:} any
reformulation in which \(\nabla u\) is not recovered by an \(L^\infty\)
Calder\'on--Zygmund bound. Two are on record: the profile reformulation, whose
status is Wall~2 --- and which \Sn{48} showed never incurs the logarithm at
all; and a symmetry-reduced setting in which \(\nabla u\) is recovered by
lower-dimensional elliptic problems, which is untested and must be settled by a
written estimate before anything is built on it.

\begin{obsv}[Wall 1 is the book's most important \OPEN]
\label{obs:wall1-open}
It would have been easy, and flattering, to classify this one \IMPOSSIBLE{}.
Three sessions had run into it independently; each concluded it could not get
past. But ``three attempts failed'' is not a proof, and \Sn{47} declined to
promote it. Chapter~\ref{ch:cz-log} is devoted to it, and to why the question
of whether it can be removed is the sharpest technical question the programme
has posed.
\end{obsv}

\section{Wall 2 --- the profile route and hypothesis (R): \texorpdfstring{\WITHDRAWN}{withdrawn}}
\label{sec:wall2}

This entry was written at \Sn{47} and classified \OPEN{}. \Sn{48} settled
\textbf{(R)} and the entry is now \WITHDRAWN{}. The document retains the \Sn{47}
text in parts (a) and (b), because the record is the point of the section, and
rewrites (c) and (d). This chapter does the same.

\textbf{Route.} Type-I blow-up-profile compactness
(Chapter~\ref{ch:profile-rigidity}): pass to a rescaled limit \(U\) and work on
the profile, where \texttt{prop:transfer-audit}(d) (line~6955) states that ``on
the profile the logarithmic factor \(\Lfac\) may be replaced by an absolute
constant''. If unconditional, Wall~1 would be removed for every estimate
transferred to \(U\).

\textbf{Obstruction, as recorded at \Sn{47}.} It was held not to be
unconditional. \texttt{prop:transfer-audit}(d) was labelled
\PROVEDMOD{(R)}, and its proof stated the mechanism plainly: ``the CZ step is
replaced by the interior estimate (R), which is logarithm-free''. Hypothesis
\textbf{(R)} \emph{is} an assumed log-free gradient bound. On that reading the
barrier was relocated, not removed.

\textbf{Classification.} \WITHDRAWN{}. The assertion that the barrier is
``relocated, not removed'' is retracted. \texttt{thm:s48-R} proves
\textbf{(R)}, and \texttt{cor:s48-transfer-d} proves
\texttt{prop:transfer-audit}(d) unconditionally. Why the \Sn{47} reading was
wrong is recorded in \texttt{rem:s48-subcritical} (line~9822): the hypothesis
actually available on the rescaled family is an \(L^\infty_tL^\infty_x\)
\emph{velocity} bound --- line~5565 assumes it outright, alongside the
vorticity bound --- and the pair \(s'=s=\infty\) satisfies \(2/s'+3/s=0<1\),
the \emph{strict}, subcritical range of Serrin's criterion, where boundedness
suffices and no smallness is required. \Sn{47} looked for \textbf{(R)} on the
\(\varepsilon\)-regularity shelf, where the smallness-versus-boundedness
objection is real; the statement is not on that shelf.
\texttt{rem:s48-mechanism} corrects the mechanism stated inside
\texttt{prop:transfer-audit}(d): the logarithm is not \emph{removed} on the
profile, it is never \emph{incurred}, because the profile route does not pass
through the Calder\'on--Zygmund operator.

\textbf{Consequence.} \emph{Permitted, and now correct:} describing the profile
route as logarithm-free without carrying \textbf{(R)}. \emph{Forbidden:}
reading this withdrawal as removing Wall~1 --- Wall~1 concerns \(u\) on
\(\T^3\) measured against \(M(t)\), and stands; the closest the result comes to
it is \texttt{obs:s48-wall1} (line~10400), recorded there as an open item and
used nowhere. \emph{Forbidden:} reading it as progress on the Prize; Wall~3 is
unaffected and \textbf{(N)}, \textbf{(N\('\))}, \textbf{(A-up)} and
\texttt{target:disorder-depletion} are exactly where \Sn{47} left them.
\emph{The one live successor:} the saturation proviso inside
\texttt{prop:transfer-audit}(d), which part~(e) uses and which \textbf{(R)}
never covered. \emph{Also forbidden, permanently:} attempting to weaken the
standing hypotheses to a vorticity-only Type-I bound ---
\texttt{prop:s48-bridge-fails} (line~10086) shows the velocity rate cannot be
recovered from \(\norm{\omega}_{L^\infty}\) and the energy, the interpolation
exponent being forced to \(3/5\), not \(1/2\).

\section{Wall 3 --- the ceiling is below the Prize}
\label{sec:wall3}

\textbf{Route.} The alignment pincer of \Sn{32}--\Sn{39}, closed at both horns.

\textbf{Obstruction.} \texttt{rem:s39-distance} (line~7703). If \textbf{(A-up)}
and \texttt{target:disorder-depletion} were proved and \textbf{(R)},
\textbf{(N)}, \textbf{(N\('\))} discharged, the conclusion would be \emph{no
Type-I singularity} --- ``strictly less than the Clay problem''. The reason is
not a gap but the construction: every estimate uses the Type-I rates, and the
profile extraction uses them \emph{essentially} --- without
\(M(t)\lesssim(T-t)^{-1}\) the rescaled family has no uniform bound and there
is no limit object. ``A Type-II singularity would rescale to nothing, and the
entire reframe would be vacuous against it.'' Two further reservations:
\textbf{(N)} restricts the conclusion to non-degenerate points, and the
surviving target must beat the constant \(1\) exactly --- the marginal pattern
the \Sn{31} audit identified as where errors hide.

\textbf{Classification.} \OPEN{}, and outside the present methods.

\textbf{Consequence.} \emph{Forbidden:} reading any part of \Sn{33}--\Sn{44} as
progress on Type-II; the document says so in terms. \emph{Forbidden:} reporting
a closed pincer as the Prize. \emph{Permitted:} closing the pincer as a Type-I
exclusion result, provided it is stated as such --- and, in a symmetry class
where Type-I is already excluded in the literature, only if the prior exclusion
is cited rather than re-derived.

\begin{obsv}[Wall 3 has just been made concrete]
\label{obs:wall3-concrete}
Until \Sn{49}, Wall~3 was a statement about a hypothetical. It is no longer.
The blow-up profile described in Chapter~\ref{ch:openai} has azimuthal
component \(u_\vartheta\sim\tau^{-(1/2+h)}\) with \(h>0\), which diverges
strictly faster than velocity Type-I: it is a \textbf{Type-II} profile. So the
one explicitly constructed singular geometry now in the literature for a
Navier--Stokes-type system sits precisely in the regime this programme's
machinery cannot see. That is recorded because it cuts against the programme,
and it is the sharpest available illustration of what Wall~3 excludes.
\end{obsv}

\section{Wall 4 --- level iteration for the annulus residual}
\label{sec:wall4}

\textbf{Route.} \texttt{rem:level-iteration} proposed absorbing the annulus
residual by a nested family of level cutoffs \(\chi_k\) with Young parameters
\(\delta_k=\delta_08^{-k}\), on the ``geometric beats geometric'' principle.

\textbf{Obstruction.} Three propositions in \texttt{sec:annulus-s41}, each
independent of the others: \texttt{prop:s41-\allowbreak{}nocontraction} (line~8189), the
recursion never contracts, \(\Lambda_k\ge2^84^k>1\) for every sequence
\((\delta_k)\); \texttt{prop:s41-vacuous} (line~8224), the unrolled inequality
is implied by \(Y_0\le Y_0\); \texttt{prop:s41-coarea} (line~8303), by a
co-area identity the residual produced by \emph{any} admissible cutoff is a
level density of the underlying measure, so no cleverer choice of levels exists
to be found. Chapter~\ref{ch:annulus} gives all three.

\textbf{Classification.} \IMPOSSIBLE{} for the route, and \SUPERSEDED{} as a
target: \texttt{target:band-absorption} (line~8468) is stated as ``the
replacement for \texttt{rem:level-iteration}'' and is itself \OPEN{}.

\textbf{Consequence.} \emph{Forbidden:} any further attempt to close the
annulus residual by choosing levels, cutoffs or Young parameters. What would
actually have to be proved is \texttt{rem:s41-endtarget} (line~8442), which
stresses that the absorbed term must carry \emph{the same} cutoff \(\eta\) as
the left-hand side --- ``this, not the value of any constant, is what defeats
the nested scheme'' --- and that one admissible level pair suffices, so the
difficulty is analytic and not combinatorial. This wall is what reclassified
\(\Aann\) from \emph{open-mechanical} to \emph{open}.

\section{Wall 5 --- \texorpdfstring{\(c_K\)}{c\_K} and band absorption are not the same hypothesis}
\label{sec:wall5}

\textbf{Route.} The document had described \texttt{conj:K-refined} and
\texttt{target:band-absorption} as ``of the same analytic type'', with the hope
that proving one would deliver the other, or that the apparently absolute
constants of the band statement made it the weaker target.

\textbf{Obstruction.} \texttt{sec:relation-s44} settles both hopes negatively:
disjoint supports (\texttt{rem:s44-disjoint}, line~8771); counterexamples both
ways (\texttt{prop:s44-b-fails}, line~8879; \texttt{prop:s44-a-fails},
line~8980); and the strength claim inverted
(\texttt{rem:s44-corrigendum}, line~8741). Chapter~\ref{ch:non-implication}
gives all four.

\textbf{Classification.} \IMPOSSIBLE{} for the implication, \emph{in a
precisely bounded sense}, and \WITHDRAWN{} for the strength claim. The bound on
the sense is in the document, not added here: by \texttt{rem:s44-scope}
(line~8855) an admissible test configuration need not come from a solution of
Navier--Stokes, so the two propositions do \emph{not} establish that the
conjectures are independent as statements about Navier--Stokes solutions ---
``that question is not decidable by any argument available here''. What is
proved is that no implication can be derived from the a priori budgets and
pointwise structure (A1)--(A3) in terms of which both are phrased.

\textbf{Consequence.} \emph{Forbidden:} treating a proof of either conjecture
as progress on the other. \emph{Permitted:} \texttt{conj:lc-s44} (line~9445),
the single hypothesis (LC) that covers both \(Y\)-halves. Its exact position is
easy to misstate: (LC) is \emph{verbatim} the \(c_K\) hypothesis --- same
normalisation, same \(1/\Lfac\) decay --- asserted on the \emph{larger} region
\(\{m\ge M/16\}\); by \texttt{rem:s44-lc-reading} it implies neither
\texttt{conj:K-refined} nor conversely, and is ``a genuinely new hypothesis of
an old form''. It is weaker than the alternative common hypothesis (JLE-1) in
the regime that matters, not weaker than \texttt{conj:K-refined}. Neither the
angular half nor the constant condition \(CC_1\le\delta\nu\) is covered by any
of them.

\section{Wall 6 --- the \texorpdfstring{\(c_K\)}{c\_K} numerics cannot discriminate}
\label{sec:wall6-walls}

\textbf{Route.} Measure \(c_K\) in direct numerical simulation and use the
result as evidence for or against \texttt{conj:K-refined}. Four sessions did
so: \Sn{34}, \Sn{35}, \Sn{42}--\Sn{43}, \Sn{45}.

\textbf{Obstruction.} The measurement is uninformative for a logical reason,
independent of what any run reports.

\begin{enumerate}[label=\textup{(\roman*)},leftmargin=2.4em]
\item \texttt{conj:K-refined} and \texttt{conj:lc-s44} are \emph{conditional}
  statements: each asserts a bound near a putative Type-I singularity, in the
  limit \(\Lfac\to\infty\). A solver that remains globally regular is never in
  that regime, so neither the conjecture nor its negation predicts anything
  about the numbers such a solver produces. Both predict the same thing ---
  nothing --- and a measurement that both sides of a disjunction predict
  cannot separate them.
\item \texttt{prop:regime-A-\allowbreak{}impossible} (line~6514) removes the one reading
  under which a regular run could still have served as a proxy. At a genuine
  blow-up time \(\limsup_{t\to T^-}\Omega(t)=\infty\); a resolved simulation is
  a bounded-enstrophy object throughout. So the regime the conjectures speak
  about is not merely one the runs did not visit; it is not of the type a
  resolved run represents, and the decay is not something such a run could
  exhibit.
\item Independently, \texttt{rem:s44-numerics} (line~9396) states that the
  \(O(1)\) measurements of \Sn{34}--\Sn{35} are ``consistent with, but do not
  establish'' either conjecture, that a measured \(c_K\approx1\) does
  \emph{not} support \(\Xi=O(1)\) after the \Sn{44} correction, and that the
  discrepancy is worst exactly in the thin, filamentary configurations expected
  near a singularity.
\end{enumerate}

\textbf{Classification.} \IMPOSSIBLE{} as a test of the conjectures. The ground
is (i)--(ii), an observation about the \emph{form} of the statements; no new
mathematics is required for it, and it does not depend on resolution, viscosity
or initial condition. The diagnostic suite itself is \OPEN{} and valid for
other purposes.

\textbf{Consequence.} \emph{Forbidden:} measuring \(c_K\),
\(c_K^{\mathrm{band}}\) or the band ratios again \emph{as evidence} for or
against those conjectures; they may be computed as context. \emph{Permitted:}
measurements whose answer differs between competing mechanisms \emph{in a
regular flow} --- questions about a mechanism, not proxies for blow-up. The
\Sn{45} data are fenced in the document inside a paragraph marked ``outside
this document'', with the explicit note that they do not refute the conjecture
and that no statement in the document rests on them.
Chapter~\ref{ch:numerical-lessons} is about this wall.

\section{Wall 7 --- the \texorpdfstring{\Sn{35}}{S35} second good term does not exist}
\label{sec:wall7}

\textbf{Route.} \texttt{prop:K-self} (line~6026), \emph{provisional} as stated,
recorded an identity with \emph{two} good terms on the left: the \(K\)-density
and, additionally, the critical quadratic with a favourable weight, presented
as reinforcing the second rung.

\textbf{Obstruction.} \texttt{lem:w2-cancel} (line~6166): the \(+\nu w\ehat\)
term of the direction equation contributes exactly \(+\nu w^2\) to
\(\tfrac12\Dt w\), which cancels the sink \(-\nu\iint m^2w^2\eta^2\)
\emph{identically}.

\textbf{Classification.} \WITHDRAWN{}, and repaired. The \(K\)-control
conclusion stands with amended constants; the critical quadratic is handled
solely by \texttt{prop:second-rung}, which suffices, so no downstream
conclusion depended on the cancelled term.

\textbf{Consequence.} \emph{Forbidden:} citing the identity in its original
two-term form, or invoking the \(w^2\) term as an independent source of
smallness. ``This is the smallest of the walls and is recorded because it is
one of three instances --- with Wall~5(iii) and the corrigendum of \Sn{34} ---
of the same failure mode: a favourable statement written down before the
computation producing it was closed.'' Note also
\texttt{rem:if-verified}: if couplings C1--C2 verified in full,
\texttt{conj:K-refined} would become \emph{unnecessary}. That conditional
remains conditional.

\section{Wall 8 --- localized Constantin--Fefferman does not transfer to profiles}
\label{sec:wall8}

\textbf{Route.} Carry the H2 / Regime-B analysis of \Sn{37}, whose engine is
\texttt{prop:localized-cf} (line~6381), across the rescaling limit onto the
Type-I profile \(U\), so that the \Sn{39} reframe inherits it.

\textbf{Obstruction.} \texttt{rem:s39-nontransfer}(i) (line~7045): the
far-field step bounds
\(\int_{\abs{y}\ge\rho'}\abs{\omega(x^*-y)}\abs{y}^{-3}\dee y\) by
\(C(\rho')^{-3/2}\Omega^{1/2}\) with \(\Omega=\norm{\omega}_{L^2}^2\), and for
\(U\in\mathcal{A}(c_I)\) on \(\R^3\) there is no reason for
\(\norm{\omega_U(\cdot,s)}_{L^2(\R^3)}\) to be finite; the naive substitute
\(\abs{\omega_U}\le M_U\) makes the integral logarithmically divergent at large
\(\abs{y}\). ``The entire H2/Regime-B analysis is unavailable on the profile as
written.''

\textbf{Classification.} \IMPOSSIBLE{} \emph{as written} --- the obstruction is
a divergence, not a missing estimate. Whether some localized substitute exists
on \(\R^3\) is \OPEN{} and untried.

\textbf{Consequence.} \emph{Forbidden:} ``Any future session that uses
\texttt{prop:localized-cf} on a profile is in error.'' The companion
prohibition, \texttt{rem:s39-nontransfer}(ii), is equally binding and easier to
violate by accident: the Bridge Lemma may \emph{not} be pulled back, because
\texttt{lem:sigma-semicontinuity} gives semicontinuity in the wrong direction
and \(\sstar(U)=0\) does not give \(\sstar(u_{\lambda_j})\to0\). Any argument
that routes a profile conclusion back through H1 on the approximating sequence
is invalid; the \Sn{39} dichotomy is built to obtain its contradiction entirely
on \(U\) for exactly this reason.

\section{Wall 9 --- forced blow-up constructions cannot be de-forced}
\label{sec:wall9}

This entry is recorded \emph{outside the mathematics} of the proof document.
It carries no label there because no part of the document uses it; it is
included because it bounds what future sessions may import. Its source is the
\Sn{40} session log.

\textbf{Route.} Import the mechanism of a published finite-time blow-up
construction for the \emph{forced} Navier--Stokes equations --- an axisymmetric
anisotropic self-similar collapsing vortex core, with
\(\ell_r\asymp\tau^{1/2}\), \(\ell_z\asymp\tau^{1/2-h}\) and
\(\mathrm{Re}_\vartheta\to\infty\) while \(\mathrm{Re}_r=O(1)\) --- by deleting
the forcing term, and read the result against the Prize equations.

\textbf{Obstruction.} The forcing is load-bearing at three separate points of
that construction, not decorative: an exponentially small external force seeds
each oscillatory pulse; what remains of the momentum residual after pulse
cancellation \emph{becomes} part of \(f\); and the compact-support cutoff
introduces a further force term. The forcing is \emph{defined} as the
Navier--Stokes residual of the constructed field. Setting \(f\equiv0\) is
therefore not a simplification of the construction but a different question,
which the source does not claim to answer.

\textbf{Classification.} \IMPOSSIBLE{} as an import route --- there is nothing
to de-force, since the residual is the forcing by construction. The underlying
blow-up question for the unforced equations is of course \OPEN{}; that is the
Prize.

\textbf{Consequence.} \emph{Forbidden:} forcing terms, modified equations and
``weak internal forces'' as mechanisms anywhere in this programme;
\emph{forbidden:} reading a forced blow-up result as evidence about the
unforced equations in either direction. \emph{Permitted:} the \emph{geometry}
of such constructions as a source of adversarial initial data for the unforced
solver, which is how \texttt{layer3/aniso\_collapse\_IC.py} was built and fed
to an unmodified \(f\equiv0\) solver. The one substantive observation, recorded
as an observation and not as evidence: sustaining that anisotropic collapse
against dissipation required an externally injected, precisely tuned force.

\begin{record}
\textbf{Wall 9's framing needs one correction, made at \Sn{49}.} The wall as
written adds, after the classification, that ``in Fefferman's formulation the
result establishes alternatives (C) and (D), not (A)/(B)''. That sentence is
accurate, but it was read in \Sn{40} and \Sn{47} as though it were itself the
reason the work does not bear on this programme. It is not.

\smallskip
Alternatives (C) and (D) are Prize-eligible; Fefferman offered them
deliberately, ``to give reasonable leeway to solvers while retaining the heart
of the problem''. A correct proof of (C) wins the Millennium Prize. The reason
the work does not bear on this programme is the stronger and simpler one:
\textbf{(A) and (C) are not mutually exclusive}. Unforced global regularity
and forced breakdown can both be true, and most experts expect exactly that.
The landscape as of \Sn{49} --- two separate efforts, with different systems
and different claims --- is set out in Chapter~\ref{ch:openai},
\S\ref{sec:s49-two-strands}, which supersedes the \Sn{40} reading this wall was
built on. Wall~9's \emph{prohibition} is unaffected: the forcing still cannot
be deleted, and the geometry may still be borrowed.
\end{record}

\section{The status table}
\label{sec:walls-table}

{\footnotesize
\begin{longtable}{@{}l >{\raggedright\arraybackslash}p{0.27\textwidth} >{\raggedright\arraybackslash}p{0.30\textwidth} >{\raggedright\arraybackslash}p{0.24\textwidth}@{}}
\caption[The walls, with classifications]{The walls document's own status
table, reproduced with the \Sn{48} amendment to row~2 and no other change.
Four walls carry \IMPOSSIBLE{}, and every one of the four is
\emph{qualified} --- for a route, over a structure, as a test, or as an
import.}
\label{tab:walls}\\
\toprule
\# & Route & Classification & Primary reference \\
\midrule
\endfirsthead
\toprule
\# & Route & Classification & Primary reference \\
\midrule
\endhead
\bottomrule
\endfoot
0 & \S20--\S22 \(H^1\) Prize claim & \WITHDRAWN{} (8 grounds)
  & \texttt{sec:audit-s31} \\
1 & Sup-norm of a CZ image; crude \(K\) & \OPEN{} (structural)
  & \texttt{prop:crude-K} \\
1\('\) & \quad same log, annulus core & \OPEN
  & \texttt{rem:s41-shortfall} \\
1\(''\) & \quad same log, corrected band & \OPEN{}; needs \emph{decay}
  & \texttt{prop:s44-\allowbreak{}linfty-route} \\
2 & Profile route as log-free & \WITHDRAWN{} at \Sn{48}; (R) proved
  & \texttt{thm:s48-R}, \texttt{cor:s48-\allowbreak{}transfer-d} \\
3 & Pincer closed \(\Rightarrow\) Prize & \OPEN{}; ceiling is Type-I only
  & \texttt{rem:s39-distance} \\
4 & Nested level iteration & \IMPOSSIBLE
  & \texttt{prop:s41-\allowbreak{}nocontraction} \\
4\('\) & \quad unrolled form & \IMPOSSIBLE{} (vacuous)
  & \texttt{prop:s41-vacuous} \\
4\(''\) & \quad any level selection & \IMPOSSIBLE{} (co-area)
  & \texttt{prop:s41-coarea} \\
4\('''\) & \quad replacement target & \SUPERSEDED{}; \OPEN
  & \texttt{target:band-\allowbreak{}absorption} \\
5 & \(c_K\Rightarrow\) band absorption
  & \IMPOSSIBLE{} over (A1)--(A3) & \texttt{prop:s44-a-fails} \\
5\('\) & band absorption \(\Rightarrow c_K\)
  & \IMPOSSIBLE{} over (A1)--(A3) & \texttt{prop:s44-b-fails} \\
5\(''\) & \quad scope of 5, 5\('\)
  & Test configurations, \emph{not} NS solutions & \texttt{rem:s44-scope} \\
5\('''\) & band is the weaker target & \WITHDRAWN{} (lemma inverted)
  & \texttt{rem:s44-corrigendum} \\
5\('''\,'\) & common hypothesis (LC) & \OPEN{}; covers both \(Y\)-halves
  & \texttt{conj:lc-s44} \\
6 & \(c_K\) numerics as evidence & \IMPOSSIBLE{} (logical)
  & \texttt{prop:regime-A-\allowbreak{}impossible} \\
6\('\) & \quad in-document assessment & Consistent with, does not establish
  & \texttt{rem:s44-numerics} \\
7 & \Sn{35} second good term & \WITHDRAWN{} (exact cancellation)
  & \texttt{lem:w2-cancel} \\
8 & Localized CF on profiles & \IMPOSSIBLE{} (\(\Omega=\infty\))
  & \texttt{rem:s39-\allowbreak{}nontransfer}(i) \\
8\('\) & Pull-back of H1 to the sequence
  & Forbidden (no reverse inequality) & \texttt{rem:s39-\allowbreak{}nontransfer}(ii) \\
9 & De-forcing a forced blow-up & \IMPOSSIBLE{} as an import
  & \Sn{40} log (outside the document) \\
\end{longtable}}

\section{Two remarks the document closes with}
\label{sec:walls-closing}

\subsection{The handover corrigenda}

\texttt{rem:s47-handover-corrigenda} records that the execution plan which
commissioned the walls document states two things the document does not
support, and that the document governs.

\begin{enumerate}[label=(\roman*),leftmargin=2.4em]
\item It refers to ``the S31 audit (all seven grounds)''.
  \texttt{sec:audit-s31} lists \textbf{eight}, D1--D8. D8 is the one most
  easily dropped from a summary and is also the first appearance of Wall~1, so
  the omission is not neutral.
\item It describes \texttt{conj:lc-s44} as ``a common \emph{weaker} target''.
  By \texttt{rem:s44-lc-reading} it neither implies nor is implied by
  \texttt{conj:K-refined}.
\end{enumerate}

A third statement is correct but routinely rounded off: the \Sn{44}
non-implications are proved over (A1)--(A3), \emph{not} for Navier--Stokes
solutions. ``\emph{Proved independent} without that qualifier overstates
them.''

\subsection{What the section does not do}

\begin{record}
\texttt{rem:s47-scope}, quoted. ``It does not assess whether any wall can be
broken, and it contains no encouragement to try. Walls~1, 2, 3 and
5\('''\,'\) are open questions with no proof against them; Walls~4, 6, 8 and 9
carry proofs, and the proofs are in the propositions cited, not here. Nothing
in this section bears on H2, on \(\sstar\)-decay, on hypothesis \textbf{(R)},
on gap \textbf{(A-up)}, on the withdrawn material of \S20--\S22, or on any
Prize or Claim~A conclusion, all of which retain the status they had before it
was written.''

\smallskip
One item in that list has since moved: hypothesis \textbf{(R)} is now a
theorem, and Wall~2 is withdrawn in consequence. The rest stand.
\end{record}

\begin{obsv}[What the shape of the wall list says]
\label{obs:wall-shape}
Ten walls. Three are \WITHDRAWN{} claims of the programme's own --- an entire
Prize claim, a relocated barrier, a cancelled term. Four are \IMPOSSIBLE{},
each qualified, and three of those four were produced by sessions
commissioned to do the opposite. Two are \OPEN{} and structural, and one of
those --- Wall~3 --- is not a difficulty at all but a limit on what the
construction can conclude even if everything works.

\smallskip
Only one wall, Wall~1, is a live technical problem in the ordinary sense: a
thing that might yield to an idea. That is why it has a chapter of its own.
\end{obsv}
